\chapter{Dataset and Corpus Engineering}
\label{chap:dataset}

% A data chapter belongs before the methodology whenever the thesis builds its own
% corpus: the methodology can then refer to the data instead of describing it.
\Cref{sec:data-corpus} describes where the samples come from and how they are labelled, \Cref{sec:data-pipeline} the pipeline that turns them into the representation the models read, and \Cref{sec:data-stats} reports what the corpus contains.

\section{The Corpus}
\label{sec:data-corpus}

Lorem ipsum dolor sit amet, consectetur adipiscing elit, sed do eiusmod tempor incididunt ut labore et dolore magna aliqua~\cite{example_Website_2026}.
State where each sample comes from, how it was labelled, and which decisions a later chapter depends on.

\begin{margintable}
  \centering
  \caption[Corpus composition]{What the corpus holds, per class.}
  \label{tab:data-corpus}
  \small
  \begin{tabular}{@{}lr@{}}
    \toprule
    Class & Samples \\
    \midrule
    negative & 502 \\
    positive & 213 \\
    \midrule
    total & 715 \\
    \bottomrule
  \end{tabular}
\end{margintable}

Ut enim ad minim veniam, quis nostrud exercitation ullamco laboris nisi ut aliquip ex ea commodo consequat (\Cref{tab:data-corpus}).
% Source: name the query, script, or export that produced every number above.

\section{From Raw Data to the Model's Input}
\label{sec:data-pipeline}

Duis aute irure dolor in reprehenderit in voluptate velit esse cillum dolore eu fugiat nulla pariatur.

\begin{figure}[htbp]
  \centering
  \begin{tikzpicture}[x=1mm, y=1mm, font=\small,
    stage/.style={draw=figboxdraw, fill=figbox, rounded corners=1mm, minimum height=9mm, inner sep=2mm},
    arr/.style={-{Latex[length=2mm]}, figline}]
    \node[stage] (raw) at (0,0)   {raw sample};
    \node[stage] (ext) at (33,0)  {extraction};
    \node[stage] (rep) at (66,0)  {representation};
    \node[stage, draw=figaccent, fill=figaccentbg] (db) at (102,0) {database};
    \draw[arr] (raw) -- (ext);
    \draw[arr] (ext) -- (rep);
    \draw[arr] (rep) -- (db);
  \end{tikzpicture}
  \caption[The corpus pipeline]{The pipeline from a raw sample to the stored representation. The accent colour marks the artifact the later chapters read from.}
  \label{fig:data-pipeline}
\end{figure}

Excepteur sint occaecat cupidatat non proident, sunt in culpa qui officia deserunt mollit anim id est laborum (\Cref{fig:data-pipeline}).

\section{What the Corpus Contains}
\label{sec:data-stats}

Sed ut perspiciatis unde omnis iste natus error sit voluptatem accusantium doloremque laudantium, totam rem aperiam.
Report the counts a reader needs to judge the results, and give every one of them a source comment.
